\chapter*{上编\quad 辑佚}
\phantomsection
\addcontentsline{toc}{chapter}{上编\quad 辑佚}
\markboth{上编\quad 辑佚}{上编\quad 辑佚}
\chaptersubtitle{ЧАСТЬ ПЕРВАЯ}{ТЕКСТ}

\noindent
凡例列于本编之前。所收残简三十七片，分四卷：营建、岁计、撤离、后题。

\begin{center}
\small
\begin{tabular}{@{}llr@{}}
\toprule
卷次 & 内容 & 片数\\
\midrule
卷一　营建 & 选址、居住环、界、净秽山的前身 & 九\\
卷二　岁计 & 口数、粮秣、水气、工籍、定额 & 十\\
卷三　撤离 & 清点、退役、装载、口令 & 十二\\
卷四　后题 & 后世题记、伪托、讹字、圆片 & 六\\
\bottomrule
\end{tabular}
\end{center}

\vspace{4mm}

\kao{凡\quad 例}

一、释文按原件行款，一字一格，不作句读。断句用顿点，是我加的，加处一律以文意可通为准。

二、补字以方框（□）标出。凡补字，校记中说明依据；无依据者不补，径作缺字。

三、缺字不可知者作「□」，可推知而字形不全者作「△」。

四、残片编号依基金会编目，自 ФВ-1963-11 起至 ФВ-1963-47 止。本书所用简次（简一至简三十七）是我的编次，与编号不同，对照表见附录一。

五、校记只处理字，不处理意思。凡涉及释读之判断，另见下编。

六、凡引传世文献，用中华书局点校本，出注卷次与页码。

\chapter*{卷一\quad 营建}
\phantomsection
\addcontentsline{toc}{chapter}{卷一\quad 营建}
\markboth{卷一\quad 营建}{卷一\quad 营建}

\noindent
卷一九片，都不是完整的文书。它们像是从账册上剪下来的条，每条只留下一个片段。

---

\begin{jian}{简一}
□山之下，去海三舍。地无草木，风自西来，终岁不息。以石为垣，垣厚四丈，内为环，环各有门，门不可同启。
\end{jian}
\chuchu{原件残片 11 号，长 21.3 cm，金属，边缘整齐，似为刻后截断。}
\xiaojì{「□山」，上字蚀尽。「去海三舍」，三舍之制不可考，若以三十里为一舍计，则九十里，与月面所见地貌不合，疑「舍」为当地里程单位。「门不可同启」，此句独立，与前后文不相属。}

\begin{jian}{简二}
三环以上，居者稀。一环之口，常不过千。环尽处为界，界外非所宜往。
\end{jian}
\chuchu{原件残片 13 号，长 17.6 cm，陶瓷，正面有釉，背面粗糙。}
\xiaojì{「三环以上」，指居住环三号以上，见简三、简五。「界」字凡四见，简九有专述，此处只记其不可越。}

\begin{jian}{简三}
一环。二环。三环。四环。五环。六环。环各一门，门各有籍。籍在门吏，不入库。
\end{jian}
\chuchu{原件残片 15 号，长 24.1 cm，金属，字口深。}
\xiaojì{六个环名连续刻写，无间隔，与简一所称「环各有门」合。此片可视为居住环总数之确证：凡六环。}

\begin{jian}{简四}
气自下而上，水自上而下，皆以管。管在壁中，不可见。凡管，岁一检，检者不入环。
\end{jian}
\chuchu{原件残片 18 号，长 19.8 cm，金属，右端断裂。}
\xiaojì{「检者不入环」，谓检修者不进入居住环内作业，与简五「外工」之名相合。此类安排在残片中凡三见，非孤例。}

\begin{jian}{简五}
外工三百二十人，居垣外小舍。月一入环，入不过三时。出则浴，浴水不入渠。
\end{jian}
\chuchu{原件残片 21 号，长 22.4 cm，陶瓷。}
\xiaojì{「浴水不入渠」，谓其洗浴之水另作处理，不与环内水系统相接。此句与全书所记之净秽制度相关，详见下编第二章。}

\begin{jian}{简六}
垣周九里，高十二丈。垣外无墙，以地为限。限外百步，不得置物。
\end{jian}
\chuchu{原件残片 24 号，长 15.2 cm，金属，下端蚀去一角。}
\xiaojì{「垣周九里」与简一「垣厚四丈」为同一条记载之两段，中间脱去一句，脱文不可补。}

\begin{jian}{简七}
第□环西储物间，架十二，架各五层。层以板隔，板下有记。
\end{jian}
\chuchu{原件残片 28 号，长 11.7 cm，陶瓷，此片即本书三十七片之出土地点之自记。}
\xiaojì{「第□环」，环数字蚀。此片背面另有刮痕一道，非字。照片所见储物间西壁架数与此合，惟层数不可辨。}

\begin{jian}{简八}
凡物入库，先记于板，后记于册。板与册不同者，以板为正。板坏则移其字，不移其板。
\end{jian}
\chuchu{原件残片 31 号，长 18.9 cm，金属。}
\xiaojì{此条为归档通则，非某一次具体记录。板者，层板下之小牌；册者，总账。二者相异时以板为正，可见板为当时之物证，册为后录。本条与本书编次之法有直接关系，见导论第五节。}

\begin{jian}{简九}
界者，非垣非墙，非水非火。在其内者出，必失其半。失者不复。界外有山，山名不可呼。
\end{jian}
\chuchu{原件残片 34 号，长 26.8 cm，金属，为卷一最长一片。}
\xiaojì{「必失其半」，语义不明，或指所携之物，或指其人。「山名不可呼」，与简一所称「□山」相应；此山当即净秽山之前身，然本片未及其用。}
